\section{审计结论：\sieve{} 的真实增量}
\label{sec:delta}

\subsection{已被占领的主张：新颖性红线}

审计的第一个产出是一份``不能再作为主要新颖性表述''的清单。以下六类表述在 2023--2026 年文献中都有明确先例，任何一条若出现在论文 v2 的摘要、引言或贡献列表的领句位置，都会立即暴露给熟悉文献的审稿人，其效果等同于重复 IP\&MMM 已经给出的否决理由。

第一，``训练无关的提示/上下文压缩''。Perception Compressor 的标题即为 A Training-Free Prompt Compression Framework（NAACL 2025 Findings），QASPC 摘要自述 training-free，NAACL 2025 综述已经把这一分支系统化成图谱。这个词组在 2026 年只能作为限定语使用，不能作为贡献。第二，``问题感知/问题条件化压缩''。LongLLMLingua 的对比困惑度、Perception Compressor 的引导问题检索、QASPC、QUITO、QGC、RECOMP 全部具备。第三，``位置感知压缩''。LongLLMLingua 的文档重排与 Perception Compressor 的 dual-slope 位置分配器是两个独立先例，且都明确以 lost-in-the-middle 为动机。第四，``句子级选择''。QASPC（外加从句级）、RECOMP 抽取式、CPC、TextRank、MMR。第五，``冗余控制''。MMR（1998）是显式惩罚项的起源，QASPC 在两层粒度上实现。第六，``自适应预算''。AdaComp 的核心贡献即自适应压缩率，LLMLingua 家族有完整的预算搜索机制。

这份红线的直接推论是：\sieve{} v1 标题所堆叠的全部定语——Training-Free、Model-Free、Question-Conditioned、Context Compression——每个词单独或两两组合都已被占位。标题必须在重构后收窄到先例未覆盖的层面（formulation 或 frontier），第 \ref{sec:retitle} 节给出候选。

\subsection{仍然站得住的差异：三根支柱}

审计确认 \sieve{} 有三根在现有文献中找不到先例的支柱，它们应该取代组件清单成为 v2 的贡献骨架。

支柱一：零模型依赖的结构性定位。全部 training-free 竞品的实际基础设施是：Perception Compressor 需要驻留 LLaMA-2-7B（压缩）加 SentenceBERT（检索）并预生成引导问题；LongLLMLingua 与 LLMLingua 需要托管小 LM 并执行多次前向；QASPC 的机制未知但属性与上述家族一致。\sieve{} 的运行时是纯 NumPy 统计，在 CPU 单核上以毫秒计（6{,}000 词上下文 53.7\,ms），与 GPT-2 门的同硬件对照是 255--320 倍的量级差。这不是渐进改进而是类别差异：免模型意味着免 GPU、免 tokenizer、免权重分发、免模型版本管理，部署面从``有一个推理服务''扩展到``有一段 Python''。审计同时给出这条支柱的防守要点：BM25 同样免模型，因此 v2 必须把主张精确表述为``在免模型象限内首次引入位置感知与冗余控制组件并给出完整三轴刻画''，而非``免模型压缩''本身——后者会被 BM25 一句反杀。

支柱二：显式的三目标联合选择形式化。MMR 形式化了相关性--冗余二元权衡，LongLLMLingua 与 Perception Compressor 实现了位置机制但未形式化为选择目标，AdaComp 自适应的是预算规模而非目标结构。``在词预算下联合优化相关性覆盖、位置可靠性分配与冗余惩罚''这一问题表述在审计范围内没有先例；\sieve{} 论文 v1 的 Problem formulation 一节（含三向权衡与选择成本核算小节）实际上已经写了这个形式化，但它被埋在方法叙事之下、且没有被标题与摘要领句。v2 的工作是把它从论文第三节的正文提升为论文的身份本身。

支柱三：三轴 frontier 评估方法论与结构性发现。压缩器自身成本作为独立测量轴在审计范围内无先例（BF 列全空）；同硬件神经门复刻给出的是``该机制族的成本下界''这一可复用测量物；而 BM25 原始召回高于 \sieve{} 的负结果被正面解释为相关性目标与位置/冗余目标之间的 Pareto 权衡，配合消融（relevance-only 召回 43.6\% 对全配置 36.2\%，flat 位置先验回升至 43.3\%）与外部 BM25（43.1\%）的相互印证，构成一个``不同选择目标产生结构性 trade-off''的发现，而不仅是``我们的方法比基线好''。这条支柱的额外价值在于它是可扩展的：后续任何压缩器（包括 QASPC 若公开机制）都可以被摆上同一条 frontier 定价，论文由单点方法变成测量框架。

\subsection{组合新颖性的残余风险与辩护策略}

三根支柱并不能完全消除审稿风险，残余风险集中在两点，v2 需要预先布防。第一点是``左上象限空缺可能只是无人感兴趣而非无人做过''：免模型象限的组件贫乏或许反映学界共识——统计信号做不了位置与冗余。\sieve{} 的辩护材料恰好是论文已有的实验：位置先验的消融（全配置 36.2\% 对 flat 先验 43.3\%）与跨数据集迁移证明统计组件确实在做可测量的工作，而不是摆设；同时其局限（词汇覆盖条件、语义改写鲁棒性弱）在 v1 的形式化分析中已有诚实边界。辩护句式应当是``这个角落空缺是因为它需要回答一个别人没有问的问题：多目标选择在零基础设施约束下还剩多少可用性''，而不是``没人做过所以新''。

第二点是三根支柱之间的依赖：支柱一（零模型）的对手是 BM25，支柱二（formulation）的对手是 MMR，支柱三（frontier）的对手是无——三者单独看都面临``和 X 比新在哪''的追问，合起来才构成``用零基础设施的 instantiation 回答一个显式形式化的多目标选择问题，并以三轴方法论刻画整个设计空间''的完整闭环。因此 v2 的贡献列表必须按``形式化 $\rightarrow$ 方法论 $\rightarrow$ 实证发现 $\rightarrow$ instantiation $\rightarrow$ 可复现工件''的顺序重写，任何把 instantiation（\sieve{} 本身）放在第一条的写法都会把论文重新推回``又一个压缩器''的判决区。摘要与引言的领句同理。这套表述规则连同标题方案，在第 \ref{sec:rewrite} 节落实为逐节改写指令。
